\documentclass[12pt,a4paper]{article}
\usepackage[a4paper,margin=2cm]{geometry}
\usepackage[T2A]{fontenc}
\usepackage[utf8]{inputenc}
\usepackage[russian]{babel}
\usepackage{amsmath,amssymb,mathtools}
\usepackage{array,booktabs,longtable}
\usepackage{xcolor}
\usepackage[hidelinks]{hyperref}
\usepackage{tikz}

\setlength{\parindent}{0pt}
\setlength{\parskip}{0.65em}
\renewcommand{\arraystretch}{1.25}

\begin{document}

\begin{titlepage}
\centering
\vspace*{0.28\textheight}
{\LARGE\bfseries Ревью домашней работы\par}
\vspace{1.5cm}
{\large Автор: Лёвин Артём Александрович\par}
\vspace{0.7cm}
{\large 18 сентября 2026 года\par}
\vfill
\begin{tikzpicture}[x=1cm,y=1cm]
  \draw[line width=0.7pt] (0,0) .. controls (2,0.55) and (4,-0.55) .. (6,0);
\end{tikzpicture}
\end{titlepage}

\newpage
\section*{Сводная таблица}
\small
\setlength{\tabcolsep}{3pt}
\begin{longtable}{>{\centering\arraybackslash}p{0.09\linewidth} p{0.20\linewidth} p{0.23\linewidth} p{0.16\linewidth} p{0.18\linewidth}}
\toprule
№ задания & Тема & Суть ошибки & Ваш ответ & Правильный ответ \\
\midrule
\endfirsthead
\toprule
№ задания & Тема & Суть ошибки & Ваш ответ & Правильный ответ \\
\midrule
\endhead
\hyperref[task:7]{7} & Четырёхугольник и параллелограмм Вариньона & Решение отсутствует; не использовано условие о равенстве отрезков, соединяющих середины противоположных сторон & не указан & $126$ \\
\addlinespace
\hyperref[task:10]{10} & Параллелограмм Вариньона & Решение отсутствует; не установлена перпендикулярность диагоналей исходного четырёхугольника & не указан & $140;\ 34;\ 70$ \\
\bottomrule
\end{longtable}
\setlength{\tabcolsep}{6pt}
\normalsize

\newpage
\vspace*{0.38\textheight}
\begin{center}
{\LARGE\bfseries Разбор ошибок}
\end{center}
\vfill

\newpage
\phantomsection\label{task:7}
\section*{Задание 7}

\textbf{Текст задания.}
Отрезки, соединяющие середины противоположных сторон выпуклого четырёхугольника, равны. Диагонали исходного четырёхугольника равны $18$ и $14$. Найдите его площадь.

\textbf{Ваше решение.}
Решение не выполнено: в работе указано, что задание не решено.

\textbf{Ваш ответ.}
Не указан.

\textcolor{red}{Ошибка:} решение отсутствует, поэтому условие о равных отрезках, соединяющих середины противоположных сторон, не было переведено в геометрическое свойство исходного четырёхугольника.

\textbf{Где именно возникает ошибка.}
Последнего правильного математического шага нет, поскольку решение не начиналось. Первый необходимый шаг --- обозначить середины сторон и построить параллелограмм Вариньона. Именно в нём данные из условия становятся удобными для использования.

\textbf{Почему это неверно.}
Само по себе равенство двух отрезков, соединяющих середины противоположных сторон четырёхугольника, ещё не даёт площадь напрямую. Нужно заметить, что эти два отрезка являются диагоналями параллелограмма, составленного из середин сторон исходного четырёхугольника. По теореме Вариньона середины сторон любого выпуклого четырёхугольника образуют параллелограмм. Если диагонали параллелограмма равны, то этот параллелограмм является прямоугольником. Тогда его соседние стороны перпендикулярны. Эти стороны параллельны диагоналям исходного четырёхугольника, следовательно, диагонали исходного четырёхугольника также перпендикулярны. После этого площадь находится однозначно.

\textcolor{blue!60!black}{Ключевая идея:} перенести условие о серединах сторон на параллелограмм Вариньона. Равные диагонали этого параллелограмма превращают его в прямоугольник, а значит, диагонали исходного четырёхугольника перпендикулярны.

\textbf{Исправление от последнего правильного шага.}
Поскольку математических шагов в Вашем решении нет, начинаем с первого необходимого построения. Обозначим через $K,L,M,N$ середины сторон $AB,BC,CD,DA$ соответственно. Тогда $KLMN$ --- параллелограмм Вариньона. Отрезки, соединяющие середины противоположных сторон исходного четырёхугольника, --- это $KM$ и $LN$, то есть диагонали параллелограмма $KLMN$. По условию $KM=LN$, поэтому $KLMN$ --- прямоугольник.

\textbf{Полное правильное решение.}
Рассмотрим выпуклый четырёхугольник $ABCD$. Пусть $K,L,M,N$ --- середины сторон $AB,BC,CD,DA$.

В треугольнике $ABC$ отрезок $KL$ соединяет середины сторон $AB$ и $BC$. По свойству средней линии треугольника
\[
KL\parallel AC,\qquad KL=\frac12 AC=9.
\]

В треугольнике $BCD$ отрезок $LM$ соединяет середины сторон $BC$ и $CD$, поэтому
\[
LM\parallel BD,\qquad LM=\frac12 BD=7.
\]

Аналогично $MN\parallel AC$ и $NK\parallel BD$, следовательно, $KLMN$ --- параллелограмм.

По условию отрезки, соединяющие середины противоположных сторон исходного четырёхугольника, равны. Это отрезки $KM$ и $LN$, то есть диагонали параллелограмма $KLMN$. Значит,
\[
KM=LN.
\]

Параллелограмм с равными диагоналями является прямоугольником. Поэтому
\[
KL\perp LM.
\]

Но $KL\parallel AC$, а $LM\parallel BD$. Следовательно,
\[
AC\perp BD.
\]

Пусть диагонали $AC$ и $BD$ пересекаются в точке $O$. Тогда четырёхугольник разбивается на четыре прямоугольных треугольника. Сумма их площадей равна
\[
\begin{aligned}
S_{ABCD}
&=\frac12 AO\cdot BO+\frac12 BO\cdot CO
 +\frac12 CO\cdot DO+\frac12 DO\cdot AO\\
&=\frac12(AO+CO)(BO+DO)\\
&=\frac12\,AC\cdot BD.
\end{aligned}
\]

Подставляем длины диагоналей:
\[
S_{ABCD}=\frac12\cdot18\cdot14=126.
\]

\textbf{Проверка.}
Параллелограмм Вариньона оказался прямоугольником со сторонами $9$ и $7$, поэтому его площадь равна
\[
S_{KLMN}=9\cdot7=63.
\]
По теореме Вариньона площадь параллелограмма, образованного серединами сторон четырёхугольника, равна половине площади исходного четырёхугольника. Поэтому
\[
S_{ABCD}=2\cdot63=126,
\]
что подтверждает найденный результат.

\textcolor{teal!60!black}{Итог:} площадь исходного четырёхугольника равна $126$.

\textbf{Задача на закрепление.}
В выпуклом четырёхугольнике отрезки, соединяющие середины противоположных сторон, равны. Его диагонали равны $24$ и $10$. Найдите площадь четырёхугольника.

\newpage
\phantomsection\label{task:10}
\section*{Задание 10}

\textbf{Текст задания.}
В выпуклом четырёхугольнике $ABCD$ диагонали равны $20$ и $14$. Точки $K,L,M,N$ являются серединами сторон $AB,BC,CD,DA$. Диагонали параллелограмма $KLMN$ равны. Найдите: площадь четырёхугольника $ABCD$; периметр параллелограмма $KLMN$; площадь параллелограмма $KLMN$.

\textbf{Ваше решение.}
Решение не выполнено: в работе указано, что задание не решено.

\textbf{Ваш ответ.}
Не указан.

\textcolor{red}{Ошибка:} решение отсутствует; не использованы свойства параллелограмма Вариньона и признак прямоугольника по равенству диагоналей.

\textbf{Где именно возникает ошибка.}
Последнего правильного математического шага нет, поскольку решение не начиналось. Первый необходимый шаг --- выразить стороны параллелограмма $KLMN$ через диагонали исходного четырёхугольника и использовать условие о равенстве диагоналей $KLMN$.

\textbf{Почему это неверно.}
Точки $K,L,M,N$ даны не случайно. Соединение середин последовательных сторон любого четырёхугольника образует параллелограмм Вариньона. Его стороны параллельны диагоналям исходного четырёхугольника и равны половинам этих диагоналей. Поэтому здесь сразу можно получить длины соседних сторон $KLMN$. Дополнительное условие о равенстве диагоналей $KLMN$ означает, что этот параллелограмм является прямоугольником. Отсюда следует перпендикулярность его соседних сторон, а значит, и перпендикулярность диагоналей $AC$ и $BD$ исходного четырёхугольника. После этого все три требуемые величины находятся стандартными школьными способами.

\textcolor{blue!60!black}{Ключевая идея:} параллелограмм $KLMN$ имеет стороны, равные половинам диагоналей $ABCD$. Равенство диагоналей $KLMN$ делает его прямоугольником, поэтому диагонали исходного четырёхугольника перпендикулярны.

\textbf{Исправление от последнего правильного шага.}
Поскольку решение не начиналось, сразу применим свойства средних линий. В треугольнике $ABC$ отрезок $KL$ является средней линией, поэтому $KL\parallel AC$ и $KL=\frac12 AC$. В треугольнике $BCD$ отрезок $LM$ является средней линией, поэтому $LM\parallel BD$ и $LM=\frac12 BD$. Далее используем равенство диагоналей параллелограмма $KLMN$.

\textbf{Полное правильное решение.}
Так как $K$ и $L$ --- середины сторон $AB$ и $BC$, то в треугольнике $ABC$
\[
KL\parallel AC,\qquad KL=\frac12 AC=\frac12\cdot20=10.
\]

Так как $L$ и $M$ --- середины сторон $BC$ и $CD$, то в треугольнике $BCD$
\[
LM\parallel BD,\qquad LM=\frac12 BD=\frac12\cdot14=7.
\]

Аналогично получаем $MN\parallel AC$ и $NK\parallel BD$. Поэтому $KLMN$ --- параллелограмм Вариньона.

По условию его диагонали равны. Параллелограмм с равными диагоналями является прямоугольником. Следовательно,
\[
KL\perp LM.
\]

Поскольку $KL\parallel AC$ и $LM\parallel BD$, получаем
\[
AC\perp BD.
\]

\textbf{1. Площадь четырёхугольника $ABCD$.}
При перпендикулярных диагоналях площадь выпуклого четырёхугольника равна половине произведения длин диагоналей:
\[
S_{ABCD}=\frac12\cdot20\cdot14=140.
\]

\textbf{2. Периметр параллелограмма $KLMN$.}
Его соседние стороны равны $10$ и $7$, поэтому
\[
P_{KLMN}=2(10+7)=34.
\]

\textbf{3. Площадь параллелограмма $KLMN$.}
Так как $KLMN$ --- прямоугольник,
\[
S_{KLMN}=10\cdot7=70.
\]

\textbf{Проверка.}
По теореме Вариньона площадь параллелограмма, соединяющего середины сторон четырёхугольника, равна половине площади исходного четырёхугольника. Половина от $140$ равна
\[
\frac12\cdot140=70,
\]
что совпадает с площадью прямоугольника $KLMN$, найденной через его стороны $10$ и $7$.

\textcolor{teal!60!black}{Итог:} площадь четырёхугольника $ABCD$ равна $140$, периметр параллелограмма $KLMN$ равен $34$, его площадь равна $70$.

\textbf{Задача на закрепление.}
В выпуклом четырёхугольнике $ABCD$ диагонали равны $18$ и $12$. Точки $K,L,M,N$ являются серединами сторон $AB,BC,CD,DA$, а диагонали параллелограмма $KLMN$ равны. Найдите площадь четырёхугольника $ABCD$, периметр параллелограмма $KLMN$ и его площадь.

\vfill
\begin{center}
\small Лёвин Артём Александрович эксклюзивно для Софии Кальней
\end{center}

\end{document}

% Краткое сообщение Софии: в этой работе основные решения выполнены верно. Не завершены задания семь и десять: в обоих нужно было использовать параллелограмм, образованный серединами сторон, и из равенства его диагоналей получить перпендикулярность диагоналей исходного четырёхугольника. В ревью подробно разобран именно этот ход.
